In countries where official results are published only at an electoral-district
level that does not correspond to any standard administrative unit, BLUE\_DB
codes the reporting units as electoral constituencies.
These entries are stored in \bluepath{geo/special.csv} with
\texttt{type = Electoral Constituency} and replace LAU-level rows in the
election files.

Electoral constituencies are used in 10 countries. In four, they are the
most fine-grained reporting units for which data is available:

\begin{description}[leftmargin=1.5cm, style=sameline]
  \item[IE] 43 D\'ail constituencies and 3--4 European Parliament constituencies.
  \item[IS] 6 Icelandic constituencies (\textit{kjörd\ae mi}).
  \item[MT] 10 electoral districts for Parliamentary election and a single 
  European Parliament constituency.
  \item[SI] 88 Slovenian electoral units (\textit{volilne enote}).
\end{description}

In another six countries, electoral constituencies are used alongside standard LAUs to
report part of the results:

\begin{description}[leftmargin=1.5cm, style=sameline]
  \item[BE] Electoral cantons (\textit{cantons \'electoraux / kieskantonen})
            are sub-provincial units used for legislative elections.
            While only older election do not report municipality-level,
            registered voters counts are only available at the cantonal level for
            all elections to date.
  \item[EE] Two special aggregate units. Additionally, 15 counties (\textit{maakond})
            are used in early elections for reporting registered voters and turnout.
  \item[EL] One special unit is used to report votes that cannot be attributed to a municipality.
  \item[HR] Three Croatian special units covering diaspora and domestic
            aggregates.
  \item[LT] One national aggregate used where a municipal breakdown is
            unavailable.
  \item[PT] Two units used for a grouping of two municipalities and votes
            not attributable to any municipality, respectively.
\end{description}
